\documentclass[11pt]{amsart}
\usepackage{amsmath,amssymb,amsthm}

\newtheorem{theorem}{Theorem}[section]
\newtheorem{conjecture}[theorem]{Conjecture}
\newtheorem{lemma}[theorem]{Lemma}
\newtheorem{proposition}[theorem]{Proposition}
\newtheorem{corollary}[theorem]{Corollary}
\numberwithin{equation}{section}
\newcommand{\K}{\mathcal K}
\newcommand{\one}{\mathbf 1}
\newcommand{\E}{\mathcal E}
\DeclareMathOperator{\vol}{vol}

\title[Threshold graphs maximizing Kemeny's constant]{Threshold graphs maximizing Kemeny's constant}
% Insert the final author information here, one block per author.
% \author{Author name}
% \address{Department, institution, and postal address}
% \email{Email address}
\subjclass[2020]{Primary 05C81; Secondary 05C50, 60J10}
\keywords{Kemeny's constant, threshold graph, pineapple graph, multiplicative degree-Kirchhoff index}
\date{}

\begin{document}
\begin{abstract}
We prove a conjecture of Breen et al. on Kemeny's constant for the
simple random walk on a connected threshold graph. At each order $n$,
every maximizer is a pineapple, that is, a clique with all remaining
vertices attached as leaves to one clique vertex. The maximizer is
unique up to isomorphism except at $n=21$. For $n\ge247$ its clique
has size $1+\lfloor\sqrt{2n-9}\rfloor$, and the maximum is
$n+\sqrt{n/2}-7/2+O(n^{-1/2})$ as $n\to\infty$.
At fixed order we also show that the pineapples are
precisely the vertices of the upper convex hull of the points
representing the edge count and the multiplicative degree-Kirchhoff
index. The proof uses the nesting of the neighborhoods outside a
maximum clique. We also characterize the sequences of threshold
graphs for which Kemeny's constant is $n+(1+o(1))\sqrt{n/2}$.
\end{abstract}
\maketitle

\section{Introduction}\label{sec:intro}

Let $G$ be a finite, simple, connected graph of order $n\ge2$ and
size $m$. Consider the simple random walk on $G$. Write $d_v$ for the
degree of $v$ and $H_y$ for the first hitting time of $y$, with $H_y=0$
when the walk starts at $y$. Kemeny and Snell \cite{KS} showed that
the mean $\K(G)=\sum_{y\in V(G)}d_y\mathbb E_x[H_y]/(2m)$ is
independent of the starting vertex $x$. Thus $\K(G)$ is the mean time
to hit a vertex drawn from the stationary distribution. If
$0=\nu_1<\nu_2\le\cdots\le\nu_n$ are the eigenvalues of the
normalized Laplacian $I-D^{-1/2}AD^{-1/2}$, where $A$ and $D$ are the
adjacency and degree matrices of $G$, then
$\K(G)=\sum_{j=2}^n\nu_j^{-1}$.

The form of an extremal graph depends on the class under
consideration. For instance, on trees of fixed order Kemeny's constant
is an increasing affine function of the Wiener index, the sum of the
distances over unordered pairs of vertices \cite{JKS}. We study
\emph{threshold graphs}, which are obtained from a single vertex by
successively adding either an isolated vertex or a vertex adjacent to
all earlier vertices \cite{MP}. Breen et al.\ \cite{Breen} expressed
Kemeny's constant in terms of this construction and proved that
$\K(G)<2n-3$ and $\K(G)<n-1+\tfrac32\sqrt m$ for connected
threshold graphs of order $n\ge3$. Their computations led to a
conjecture about the graphs that maximize $\K$ at each order.

For $2\le k\le n$, let $P_{n,k}$ be the graph obtained from a
$k$-clique by attaching $n-k$ leaves to a single clique vertex. We call
$P_{n,k}$ a \emph{pineapple}, including the star $P_{n,2}$ and the
complete graph $P_{n,n}$. In citing numbered statements of Breen et al.,
we use the numbering of the preprint listed in \cite{Breen}.

\begin{conjecture}[Breen et al.\ {\cite[Conjecture 3.13]{Breen}}]\label{conj:pineapple}
Among the connected threshold graphs of a given order, every graph
with maximum Kemeny's constant is a pineapple.
\end{conjecture}

We prove Conjecture~\ref{conj:pineapple} by comparing graphs at a
prescribed edge count. The quantity $2m\K(G)$ is the multiplicative
degree-Kirchhoff index of $G$. It equals
$\sum_{\{u,v\}}d_ud_v R(u,v)$, where the sum is over unordered
pairs of distinct vertices and $R(u,v)$ is their effective resistance
when each edge has unit resistance \cite{CZ}. For fixed $n$, abbreviate
$P_{n,t}$ to $P_t$ and write $m_t=\binom t2+n-t$ for its edge count.

\begin{theorem}\label{th:comparison}
Let $G$ be a connected threshold graph of order $n\ge3$ and size $m$.
Choose $2\le t<n$ such that $m_t\le m\le m_{t+1}$, and set
$\lambda=(m-m_t)/(m_{t+1}-m_t)$. Then
\[
 m\K(G)\le (1-\lambda)m_t\K(P_t)
               +\lambda m_{t+1}\K(P_{t+1}).
\]
Equality holds if and only if $G$ is isomorphic to $P_t$ or $P_{t+1}$.
In particular, for each $2\le t\le n$, the graph $P_t$ is the unique
maximizer up to isomorphism at order $n$ and size $m_t$.
\end{theorem}

Join consecutive pineapple points in the $(m,2m\K)$-plane.
The slopes of the resulting segments strictly decrease, and
Theorem~\ref{th:comparison} identifies their union as the upper boundary
of the convex hull of the points $(m,2m\K(G))$. Dividing the
inequality by $m$ bounds $\K(G)$ by a weighted mean of the two
pineapple values. It therefore suffices to determine which pineapple
maximizes $\K$.

\begin{theorem}\label{th:main}
Every connected threshold graph maximizing Kemeny's constant at a fixed
order $n\ge2$ is a pineapple. The maximizer is unique up to isomorphism
except at $n=21$, where the maximizing graphs are $P_{21,6}$ and
$P_{21,7}$. For $n\ge247$ the extremal clique size is
$1+\lfloor\sqrt{2n-9}\rfloor$. As $n\to\infty$, the maximum is
$n+\sqrt{n/2}-7/2+O(n^{-1/2})$.
\end{theorem}

Breen et al.\ also conjectured that the maximum is
$n+\sqrt n/2+O(1)$ and proposed two candidates for the length $r$ of
the first block of ones in the construction code
$0\mathbf 1^r\mathbf 0^{n-r-2}1$
\cite[Proposition 3.15 and Conjecture 3.16]{Breen}.
Here $0$ denotes an isolated addition and $1$ a dominating addition,
with the initial vertex encoded by $0$. The clique has size $r+2$.
Theorem~\ref{th:main} gives the coefficient $1/\sqrt2$ of $\sqrt n$.
At $n=21$ the optimal values are $r=4,5$, both giving
$\K=21$. The two candidates $r=6,7$ in the preprint give
$3429/164$ and $1495/72$, respectively.

The proof uses the nested neighborhoods outside a maximum clique.
Section~\ref{sec:averaging} expresses $\K(G)$ in terms of their
sizes and averages over a common level in the nesting. This reduces
Theorem~\ref{th:comparison} to graphs with just two possible
neighborhood sizes. In Section~\ref{sec:polygon} we compare each
such graph with a chord joining two pineapple points.
Section~\ref{sec:cliquesize} determines the optimal clique size at
every order. Section~\ref{sec:stability} characterizes the graphs
that attain the leading term of the excess $\K(G)-n$.

\section{Averaging nested neighborhoods}\label{sec:averaging}

\subsection{A formula in terms of neighborhood sizes}
Fix a construction of a connected threshold graph $G$ of order $n\ge2$.
The initial vertex and the dominating additions form a clique $C$.
Number its vertices $1,\ldots,k$ in reverse order of addition, so
that the initial vertex has label $k$. The remaining vertices form
an independent set $I$. Each $i\in I$ is adjacent precisely to the
clique vertices added after $i$. Its neighborhood is therefore an
initial segment $[r_i]=\{1,\ldots,r_i\}$, where $1\le r_i\le k-1$.
A clique meeting $I$ contains at most one vertex of $I$ and at most
$k-1$ of its neighbors, so $k=\omega(G)$.

Write $N=n-k$, $T=\sum_{i\in I}r_i$, and let
$V=\vol(G)=2|E(G)|=k(k-1)+2T$. In the bipartite adjacency matrix
with rows indexed by $I$ and columns by $C$, row $i$ has $r_i$
initial ones. For $2\le j\le k$, let $N_j$ be the number of rows
reaching column $j$, and let $A_j$ be the total length of the
remaining rows. With $f(x)=2x(1-2x/V)$, write
\[
 N_j=|\{i:r_i\ge j\}|,\qquad A_j=\sum_{r_i<j}r_i,
 \qquad w_j=\frac1{j(j-1)}.
\]

\begin{lemma}\label{lem:spectral}
For a connected threshold graph $G$ with the preceding parameters,
\begin{equation}\label{eq:formula}
 \K(G)=n-1-\frac TV+
 \sum_{j=2}^k\frac{w_j f(A_j)-1}{k+N_j}.
\end{equation}
\end{lemma}
\begin{proof}
For real-valued functions on $V(G)$ consider the energy form and
the degree-weighted variance
\begin{align*}
 \E(g,h)&=\sum_{uv\in E(G)}(g(u)-g(v))(h(u)-h(v)),\\
 \mathcal V(g)&=\sum_v d_vg(v)^2-\biggl(\sum_v d_v g(v)\biggr)^2/V .
\end{align*}
Both forms are invariant under the addition of constants. On the
subspace of functions of degree-weighted mean zero, $\mathcal V$ is
the squared norm for the inner product
$\langle g,h\rangle_d=\sum_v d_vg(v)h(v)$. In an orthonormal
eigenbasis of the random-walk Laplacian $I-D^{-1}A$ on that subspace,
the matrices of $\E$ and $\mathcal V$ are
$\operatorname{diag}(\nu_2,\ldots,\nu_n)$ and the identity.
If $M$ and $B$ are the respective matrices in any other basis modulo
constants, $M^{-1}B$ is similar to
$\operatorname{diag}(\nu_2^{-1},\ldots,\nu_n^{-1})$. Taking the trace in an
energy-orthogonal basis $(g_i)$ therefore gives
$\K(G)=\sum_i\mathcal V(g_i)/\E(g_i,g_i)$.

For each $i\in I$, take the indicator of $\{i\}$. On $C$, let
$u_j=\sqrt{w_j}(\one_{[j-1]}-(j-1)e_j)$ for $2\le j\le k$, where
$\one_{[j-1]}$ is the indicator of $[j-1]$ and $e_j$ is the $j$th
coordinate vector. The $u_j$ form an orthonormal basis of the vectors
in $\mathbb R^k$ with coordinate sum zero. Extend $u_j$ to a function
$g_j$ on $G$ by setting $g_j(i)$ equal to the mean of $u_j$ on
$[r_i]$. This mean is $\sqrt{w_j}$ if $r_i<j$ and zero otherwise.
The indicators are mutually energy-orthogonal and are energy-orthogonal
to every $g_j$ by the averaging property. The clique
contributes $k\delta_{jh}$ to $\E(g_j,g_h)$. The edges at a vertex
of degree $r$ in $I$ contribute $\delta_{jh}$ if $j,h\le r$ and
zero otherwise, since a contrast with index greater than $r$ is
constant on $[r]$. Thus $\E(g_j,g_h)=(k+N_j)\delta_{jh}$. These
$n-1$ functions form the required basis.

The indicator of $\{i\}$ contributes $1-r_i/V$ to the trace,
so the combined contribution of the indicators is $N-T/V$. The clique
vertex at position $t$ has degree $k-1+N_t$, where $N_1=N$. Using
$\sum_{t<j}(N_t-N_j)=A_j$ and the extension values gives
$\sum_v d_vg_j(v)=2\sqrt{w_j}A_j$ and
$\sum_v d_vg_j(v)^2=k-1+N_j+2w_jA_j$. The variance identity
now gives \eqref{eq:formula}.
\end{proof}

The following consequence will be used in the proof of
Theorem~\ref{th:stability}.

\begin{corollary}\label{cor:uniform}
For every connected threshold graph of order $n\ge2$ and clique number
$k\ge2$, we have $\K(G)\le n-1+\sqrt{(1-1/k)(N-T/k)/2}$.
\end{corollary}
\begin{proof}
Put $Q_j=V-2A_j$. Then
$Q_j=k(k-1)+2\sum_{r_i\ge j}r_i\le(k+N_j)^2$. Since
$xy/(x+y)\le\sqrt{xy}/2$, taking $x=2A_j$ and $y=Q_j$ gives
$f(A_j)/(k+N_j)\le\sqrt{A_j/2}$. Discarding the negative terms
in \eqref{eq:formula} and applying Cauchy--Schwarz bounds
$\K(G)-n+1$ by $(\sum_j w_j)^{1/2}(\sum_j w_jA_j/2)^{1/2}$.
Now $\sum_{j=2}^k w_j=1-1/k$. A row of length $r$ contributes
$r\sum_{j=r+1}^k w_j=1-r/k$ to $\sum_{j=2}^k w_jA_j$, so the
latter sum is $N-T/k$.
\end{proof}

\subsection{A common neighborhood level}
Assume $k\ge3$, and write $q=k-1$, $a=k-2$. For each integer $0\le s\le N$, let $H_s$ consist of the clique $C$ and an independent set
of size $N$, with $s$ vertices adjacent to $[q]$ and the remaining
$N-s$ adjacent to $[1]$. Thus $H_0=P_k$. For
a graph $X$ of order $n$, put $Y(X)=\vol(X)(\K(X)-n+4)$.
This change of ordinate subtracts the linear function $(n-4)\vol(X)$
from the multiplicative degree-Kirchhoff index and preserves every
comparison with a chord.

Choose a common level $j$ uniformly from $2,\ldots,q$. Replace
$[r_i]$ by $[q]$ if $r_i\ge j$ and by $[1]$ otherwise. The resulting
graph is $H_{N_j}$, and the new degree of each $i\in I$ has mean $r_i$.

\begin{lemma}[Common-level averaging]\label{lem:average}
With the notation above,
\[
 \vol(G)=\frac1a\sum_{j=2}^q \vol(H_{N_j}),\qquad
 Y(G)\le\frac1a\sum_{j=2}^q Y(H_{N_j}).
\]
Equality in the second relation holds precisely when all $r_i$ lie in
$\{1,q\}$.
\end{lemma}
\begin{proof}
Write $b_j=N_j$, $R_j=w_jA_j$ and $Q_j=V-2A_j$ for $2\le j\le q$.
All sums in this proof are over indices in $\{2,\ldots,q\}$ unless
otherwise indicated.
Separating the $j=k$ term in \eqref{eq:formula} gives
\begin{equation}\label{eq:affine}
 Y(G)=B(T)+\sum_j\frac{2R_jQ_j-V}{k+b_j},\qquad
 B(T)=3V-\frac{aT+V}{k}.
\end{equation}
Here $B$ is affine since $V=kq+2T$.

Two counts of the row lengths give
\[
 Q_j=j(k+2b_j)+\sum_{\ell>j}(k+2b_\ell),\qquad
 N-b_\ell=\ell R_\ell+\sum_{j<\ell}R_j.
\]
The first counts each row of length $r\ge j$ exactly $r$ times in
$T-A_j$. In the second, a row of length $r<\ell$ contributes
$r/(\ell-1)+r\sum_{j=r+1}^{\ell-1}w_j=1$, and rows of length at least $\ell$ contribute zero.
Interchanging the two sums and using $b_j\ge b_\ell$ for $j<\ell$ gives
\begin{align*}
 \sum_j\frac{R_jQ_j}{k+b_j}
 &=\sum_\ell(k+2b_\ell)\left(\frac{\ell R_\ell}{k+b_\ell}
                   +\sum_{j<\ell}\frac{R_j}{k+b_j}\right)\\
 &\le\sum_\ell\frac{(k+2b_\ell)(N-b_\ell)}{k+b_\ell}.
\end{align*}
For $H_s$ the columns $2,\ldots,q$ all have height $s$, so equality
holds throughout.

Let $p=a^{-1}\sum_j b_j=(T-N)/a$ and put $v_s=\vol(H_s)$.
Since each row keeps its mean length, the mean of $T$, and hence the
means of $V$ and $B(T)$, are preserved.
Insert the last bound in \eqref{eq:affine} for $G$, and use equality
there for each $H_{b_j}$. The terms involving $N-b_j$ cancel, leaving
\[
 \frac1a\sum_jY(H_{b_j})-Y(G)
 \ge \sum_j\frac{V-v_{b_j}}{k+b_j}
 =\frac{2a}{k+p}\sum_j\frac{(b_j-p)^2}{k+b_j}.
\]
The last equality uses $V-v_{b_j}=2a(p-b_j)$ and $\sum_j(p-b_j)=0$.
If the heights $b_2,\ldots,b_q$ are not all equal, the right side is positive.
Otherwise every row has length $1$ or $q$, so every outcome is $G$.
\end{proof}

\section{The upper convex hull}\label{sec:polygon}
Set $V_t=2m_t=t^2-3t+2n$ and $\phi(t)=2nt-n-5+4n/t$.
The value of Kemeny's constant on a pineapple is
\begin{equation}\label{eq:pine}
 \K(P_t)=n-4+\frac{\phi(t)}{V_t}
         =n-4+\frac2t+\frac{(n-1)(2t-1)}{V_t}.
\end{equation}
This is \cite[Lemma 3.14]{Breen} with $r=t-2$. A resistance-distance
derivation appears in \cite[Corollary 3.5]{HC}. Formula
\eqref{eq:formula} gives the same expression by taking $r_i=1$ throughout.

Let $F$ be the piecewise linear function joining the points
$(V_t,\phi(t))$, $2\le t\le n$. Its slope between consecutive
points is $n(t+2)/[t(t+1)]$, which decreases strictly. Thus $F$ is
concave, with a vertex at every pineapple volume.
For $k\ge3$, the graph $H_s$ has $T=N+as$ and
$\vol(H_s)=V_k+2as$. For $2\le j\le k-1$, its parameters are
$N_j=s$ and $A_j=N-s$, while $N_k=0$ and $A_k=T$.
Substitution in \eqref{eq:affine} gives
\begin{equation}\label{eq:line}
 Y(H_s)=\chi(s)=\phi(k)+2as\,\tau(s),\qquad
 \tau(s)=\frac{2n(k+1)+k(s+1)}{2k(k+s)} .
\end{equation}

For $1\le s\le N$ integral, denote the slope of the chord joining
the pineapple points of clique sizes $k$ and $k+s$ by
\[
 S_k(s)=\frac{\phi(k+s)-\phi(k)}{V_{k+s}-V_k}
       =\frac{2n}{2k+s-3}\left(1-\frac{2}{k(k+s)}\right).
\]

\begin{proof}[Proof of Theorem~\ref{th:comparison}]
For each $s\ge1$, the volume of $H_s$ lies strictly between $V_k$
and $V_{k+s}$, since $\vol(H_s)-V_k=2as$ and
$V_{k+s}-\vol(H_s)=s(s+1)$. By \eqref{eq:line}, the slope from the
point of $P_k$ to the point of $H_s$ is $\tau(s)$. The chord from
$P_k$ to $P_{k+s}$ has larger slope, since
\[
 S_k(s)-\tau(s)=
 \frac{(s+1)[2n(k-1)-k(2k+s-3)]}{2k(k+s)(2k+s-3)}>0 .
\]
Indeed $n\ge k+s$ and the bracket is at least $k+(k-2)s$.
Consequently $Y(H_s)<F(\vol(H_s))$ for $s>0$, while $H_0=P_k$
gives equality. A nonpineapple has $k\ge3$ and $N_2>0$.
Lemma~\ref{lem:average} and concavity now give
\[
 Y(G)\le\frac1a\sum_{j=2}^q Y(H_{N_j})
 <\frac1a\sum_{j=2}^q F(\vol(H_{N_j}))\le F(\vol(G)).
\]
Every pineapple attains its own vertex of $F$. Adding back the
linear term $(n-4)V$ and dividing by two proves the theorem.
\end{proof}

We next keep both the order $n$ and the clique size $k\ge3$ fixed.
Extend $\chi$ to the real interval $[0,N]$ by \eqref{eq:line}. The
identity
\[
 \chi(s)=\phi(k)+as+
 \frac{a[2n(k+1)-k(k-1)]}{k}\,\frac{s}{k+s}
\]
shows that $\chi$ is strictly concave when $N>0$. It also shows
that $\tau(s)=(\chi(s)-\chi(0))/(2as)$ decreases strictly for $s>0$,
with $\tau(0)$ given by \eqref{eq:line}, since
$2n(k+1)-k(k-1)>0$.
Let $\widehat\chi$ interpolate the integer values of $\chi$ linearly. Then
Lemma~\ref{lem:average} and $\sum_{j=2}^q N_j=ap$ give
\begin{equation}\label{eq:int}
 Y(G)\le\widehat\chi(p)\le\chi(p),\qquad p=(T-N)/a.
\end{equation}
If $p$ is integral, $H_p$ is therefore the unique maximizer up to
isomorphism at fixed order, clique size, and edge count, by the equality condition in Lemma~\ref{lem:average}.

\begin{corollary}\label{cor:clique}
Let $G$ be a connected threshold graph of order $n$ and clique number
$k$ that is not a pineapple. Let $\ell$ be the least $t\in\{2,\ldots,n\}$
such that $V_t\ge\vol(G)$. Then $\K(G)<\max\{\K(P_k),\K(P_\ell)\}$.
If $k^2\ge2n$, then $\K(G)<\K(P_k)$.
\end{corollary}
\begin{proof}
For a nonpineapple $p>0$ and $V=V_k+2ap$. By \eqref{eq:int} the
slope from $(V_k,\phi(k))$ to $(V,Y(G))$ is at most $\tau(p)$.
If $2n\le k^2$, then $V_k\le k(2k-3)$. Since
$(2k-1)/(2k-3)>1+1/k$, formula \eqref{eq:pine} gives
\[
 \frac{\phi(k)}{V_k}
 >\frac2k+\frac{(n-1)(k+1)}{k^2}
 >\tau(0)>\tau(p).
\]
Thus $Y(G)/V<\phi(k)/V_k$, proving both assertions in this range.

Suppose $2n>k^2$ and put $u=\ell-k$. For $u=1$ apply
Theorem~\ref{th:comparison}. For $u\ge2$ it suffices to show
$S_k(u)>\tau(p)$, which places $(V,Y(G))$ strictly below the chord
from $P_k$ to $P_\ell$. If $p\ge u$ this follows from
$S_k(u)>\tau(u)\ge\tau(p)$. Otherwise write $j=u-1$ and $d=p-j$.
Since $V_{k+j}<V$, we have $ad>\binom{j+1}{2}\ge j$.
Thus $0<d<1$. Since $ap=T-N$ is an integer, so is $ad$.
Hence $ad\ge j+1$ and $kd=ad+2d\ge j+1+2d>p+1$.
Using $2n>k^2$ in \eqref{eq:line} now gives
\[
 \tau(p)<\frac{n}{k+p}\left(1+\frac{1+d}{k}\right)
 \le\frac{n}{k+(j-1)/2}< S_k(u).
\]
For the middle inequality, $(1+(1+d)/k)(k+(j-1)/2)\le
k+1+d+j-1=k+p$ since $1+(1+d)/k<2$.
The last inequality reduces to $(k-2)(k+u-2)>0$.
\end{proof}

\section{The extremal clique size}\label{sec:cliquesize}
It remains to maximize \eqref{eq:pine} over the integers $2\le k\le n$.
For $2\le k<n$ put
\[
 h_n(k)=(k+2)n(k^2-2n+9)-14n-5k(k+1).
\]
Direct subtraction gives
\begin{equation}\label{eq:adj}
 \K(P_{k+1})-\K(P_k)
 =\frac{-2(k-1)h_n(k)}{k(k+1)V_kV_{k+1}} .
\end{equation}
This difference admits an explicit cutoff in the order of the graph.

\begin{proposition}\label{prop:cutoff}
Let
\[
 c_k=\frac{k^2+9}{2}-\frac{12}{k+2},\qquad t_k=\lceil c_k\rceil .
\]
For $2\le k<n$, the difference $\K(P_{k+1})-\K(P_k)$ is negative
when $n<t_k$ and positive when $n\ge t_k$, with the sole exception
of $(n,k)=(21,6)$, where it is zero.
\end{proposition}
\begin{proof}
The sign can be read from the integer $D=2(k+2)(c_k-n)$, since
\[
 h_n(k)=\left(n-\frac5{k+2}\right)D-\frac{5(k-6)(k-1)}{k+2}.
\]
For $k\ge6$ and $n\ge t_k$ both terms are nonpositive, with
equality in their sum only at $(n,k)=(21,6)$. At $n=t_k-1$ we have
$D\ge1$ and $n\ge c_k-1\ge k^2/2>5(k-6)+1$, so $h_n(k)>0$.
As $h_n(k)$ is concave in $n$ and $h_1(k)=k(k-1)(k-2)\ge0$,
it is positive for every integer $1<n<t_k$.
For $k=2,3,4,5$ the cutoffs
$t_k$ are $4,7,11,16$, respectively. The pairs
$(h_{t_k-1}(k),h_{t_k}(k))$ are $(12,-6)$, $(36,-18)$,
$(60,-56)$ and $(60,-150)$. Concavity and $h_1(k)\ge0$ give
positivity up to $t_k-1$ in these four cases as well. The negative
difference from $t_k-1$ to $t_k$ and concavity give negativity for
all $n\ge t_k$. The result now follows from \eqref{eq:adj}.
\end{proof}

\begin{proof}[Proof of Theorem~\ref{th:main}]
The structural assertion follows from Theorem~\ref{th:comparison},
since division by $m$ gives a weighted mean of the two pineapple
values. At order two the only graph is $K_2$.
Since $c_{k+1}-c_k=k+1/2+12/[(k+2)(k+3)]>1$, the cutoffs $t_k$
increase strictly. By Proposition~\ref{prop:cutoff}, the largest
maximizing clique size at every order $n\ge2$ is therefore
\begin{equation}\label{eq:all-orders}
 2+\bigl|\{j:2\le j<n,\ t_j\le n\}\bigr|.
\end{equation}
The pineapple values increase up to this index and decrease thereafter,
with a single equality between consecutive values at $n=21$. At that
order $P_6$ and $P_7$ both attain the maximum.

For $k\ge23$ one has $12/(k+2)<1/2$, hence
$t_k=\lceil(k^2+9)/2\rceil$. Also $t_{22}=246$.
If $n\ge247$ and $d=\lfloor\sqrt{2n-9}\rfloor$, then $d\ge22$
and $t_d\le n<t_{d+1}$. The unique optimal clique size is therefore
$d+1$. The cutoff $247$ is sharp, since the maximizer at $n=246$ has
clique size $23$, while the formula gives $22$.

Finally write $s=\sqrt{2n}$ and $k=s+h$ at the optimum, where $h=O(1)$.
In \eqref{eq:pine}, $V_k=2s^2+(2h-3)s+O(1)$ and
$(n-1)(2k-1)=s^3+(h-1/2)s^2+O(s)$.
The quotient is $s/2+1/2+O(s^{-1})$, which gives the stated
asymptotic expansion.
\end{proof}

\section{Near-extremal graphs}\label{sec:stability}
The leading excess $\sqrt{n/2}$ in Theorem~\ref{th:main} also
determines how many vertices outside the clique may have degree
greater than one. In the decomposition of Section~\ref{sec:averaging},
write $b=N_2$ for the number of such vertices. For $k\ge3$ all
clique vertices have degree at least two, so $k+b$ is the total number
of nonpendant vertices. If $k=2$, then $G$ is a star and $b=0$.
In particular $b$ is independent of the chosen construction.

\begin{theorem}\label{th:stability}
Let $(G_n)$ be a sequence of connected threshold graphs of order $n$,
and let $k_n$ and $b_n$ denote the preceding parameters. Then
\[
 \K(G_n)-n\sim\sqrt{n/2}\quad\Longleftrightarrow\quad
 k_n\sim\sqrt{2n}\quad\text{and}\quad b_n=o(\sqrt n).
\]
\end{theorem}
\begin{proof}
Write $s=\sqrt{2n}$ and suppress the index $n$. The proof of
Corollary~\ref{cor:uniform} gives, for $2\le j\le k$,
\[
 \frac{f(A_j)}{k+N_j}\le\sqrt{A_j/2}\le\frac{n+A_j}{2s}.
\]
Since $\sum w_j\le1$ and $\sum w_jA_j=N-T/k\le n$, the weighted
sum of the right sides is at most $s/2$. By \eqref{eq:formula},
the weighted sum of the left sides is at least $\K(G)-n+1$.

Suppose $\K(G)-n\sim s/2$. The weighted sum of the gaps is $o(s)$,
so the gap for $j=2$ is $o(s)$ because $w_2=1/2$. The second
inequality has gap $(\sqrt n-\sqrt{A_2})^2/(2s)$, whence $A_2\sim n$.
Since $\sqrt{A_2/2}\sim s/2$, the ratio of the two sides in the
first inequality tends to one. Writing $Q_2=V-2A_2$, this ratio is
\[
 \frac{2\sqrt{2A_2Q_2}}{2A_2+Q_2}\,
 \frac{\sqrt{Q_2}}{k+b}.
\]
Both factors are at most one, so both tend to one. The first has the
form $2x/(1+x^2)$ with $x=\sqrt{Q_2/(2A_2)}$, which forces $x\to1$.
It follows that
$Q_2\sim 2A_2\sim 2n$ and $k+b\sim s$. Since
$Q_2\le k(k-1)+2(k-1)b\le(k+b)^2-b^2$, we have $b=o(s)$ and $k\sim s$.

Conversely, suppose $k\sim s$ and $b=o(s)$. All but $b$ of the
$N=n-k$ independent vertices have degree one, so $N\le T\le N+(k-2)b$
gives $T\sim n$ and $V\sim4n$. Uniformly for $2\le j\le k$,
$N-b\le A_j\le T$ and $N_j\le b$ give $f(A_j)\sim n$ and
$k+N_j\sim k$. Substituting in \eqref{eq:formula} and using
$\sum_j w_j=1-1/k$ yields $\K(G)-n\sim n/k\sim s/2$.
\end{proof}

\begin{thebibliography}{9}
\bibitem{Breen} J. Breen, S. Kim, A. Low Fung, A. Mann, A.~A. Parfeni
and G. Tedesco, Threshold graphs, Kemeny's constant, and related random
walk parameters, Linear Algebra Appl. \textbf{709} (2025), 284--313.
doi:10.1016/j.laa.2024.12.022. Preprint arXiv:2310.08552v1.
\bibitem{CZ} H. Chen and F. Zhang, Resistance distance and the
normalized Laplacian spectrum, Discrete Appl. Math. \textbf{155}
(2007), 654--661.
\bibitem{HC} T. Haritha and A.~V. Chithra, Resistance distance in
$k$-coalescence of certain graphs, TWMS J. Appl. Eng. Math.
\textbf{15} (2025), 1503--1519.
\bibitem{JKS} J. Jang, S. Kim and M. Song, Kemeny's constant and
Wiener index on trees, Linear Algebra Appl. \textbf{674} (2023),
230--243.
\bibitem{KS} J.~G. Kemeny and J.~L. Snell, \emph{Finite Markov
Chains}, Springer-Verlag, New York, 1976.
\bibitem{MP} N.~V.~R. Mahadev and U.~N. Peled, \emph{Threshold Graphs
and Related Topics}, Annals of Discrete Mathematics, vol.~56,
North-Holland, Amsterdam, 1995.
\end{thebibliography}
\end{document}
